\chapter*{Abstract (Deutsch)}
\addcontentsline{toc}{chapter}{Abstract (Deutsch)}

\begin{center}
    \textbf{Kurzfassung}
\end{center}

Der Internationale Flughafen Blaise Diagne (AIBD) im Senegal verzeichnet ein
Ankunftsaufkommen, das sich nicht gleichmäßig über den Betriebstag verteilt,
sondern in sogenannten Ankunftsbänken konzentriert. In diesen Zeitfenstern
treffen die Passagier- und Gepäckströme mehrerer Flüge innerhalb kurzer Zeit
auf dieselben Ressourcen wie die Grenzkontrolle, die Gepäckausgabe und der Zoll.
Unter diesen Bedingungen reicht eine rein beschreibende Betriebsüberwachung
nicht aus, da sie lediglich bereits eingetretene Zustände abbildet. Die
vorliegende Arbeit entwickelt und evaluiert daher ein ankunftszentriertes
Framework prädiktiver operativer Intelligenz, das deskriptive Analytik,
kurzfristige Prognose, ereignisdiskrete Simulation und managementorientierte
Entscheidungsunterstützung in einer durchgängigen Prozesskette verbindet.

Über die FlightRadar24-Schnittstelle wurde ein historisches
Flugbewegungsarchiv mit 50.385 Datensätzen erhoben und in einem fünfstufigen
Validierungs- und Filterverfahren auf 22.270 validierte Datensätze zu
kommerziellen Passagierankünften reduziert. Passagier- und Gepäckaufkommen
wurden anhand des Flugzeugtyps über eine feste Sitzplatzzuordnung, einen
Auslastungsfaktor von 0,82 sowie konstante Gepäckannahmen geschätzt und den
operativen Belastungsstufen zugeordnet. Nach der Konstruktion zeitlicher,
zyklischer sowie Lag- und Rolling-Merkmale wurden vier baumbasierte
Ensemble-Verfahren für zwei Aufgaben trainiert: die Regression der
Passagierbelastung des jeweils nächsten Ankunftsereignisses und die binäre
Klassifikation von Hochlastphasen. Die betrieblichen Auswirkungen wurden
mittels einer ereignisdiskreten AnyLogic-Simulation der Prozesskette
Grenzkontrolle-Gepäckausgabe-Zoll-Ausgang über sechs Szenarien und einen
Zeithorizont von zwölf Stunden bewertet; die Ergebnisse wurden in einem
fünfseitigen Power-BI-Entscheidungsunterstützungssystem zusammengeführt.

Die deskriptive Analyse identifiziert zwei wiederkehrende Ankunftsbänke
zwischen 04:00 und 06:00 Uhr sowie zwischen 19:00 und 21:00 Uhr mit
Spitzenwerten von 203,4 bzw. 210,8 geschätzten Passagieren, getrennt durch
eine Schwachlastphase mit einem Minimum von 135,7 Passagieren um 11:00 Uhr.
Da das Gepäckaufkommen über einen festen Multiplikator aus dem
Passagieraufkommen abgeleitet wird, stellt die stündliche Gepäckreihe eine
lineare Transformation der Passagierreihe dar ($r = 1{,}00$) und enthält kein
eigenständiges zeitliches Signal; beide Reihen weisen damit dieselben
Ankunftsbänke auf.

LightGBM erzielte die beste Regressionsleistung
(MAE: 34,06; RMSE: 46,61; $R^2$: 0,2366). Bei der Stauklassifikation wurde
eine zunächst nahezu perfekte Modellgüte infolge von Informationsleckage
durch eine zukunftsbezogene Prädiktorvariable diagnostiziert; nach deren
Entfernung erreichte das eingesetzte LightGBM-Modell eine Genauigkeit von
0,7805, einen F1-Wert von 0,4690 und einen ROC-AUC-Wert von 0,7914.

Die Simulation weist die Gepäckausgabe als dominanten Engpass des
Ausgangszustands aus: Eine alleinige Erhöhung der Gepäckkapazität senkte die
durchschnittliche Verweildauer von 102 auf 83 Minuten und steigerte den
Durchsatz von 484 auf 573 abgefertigte Passagiere bei zugleich höchstem
Resilienzwert (0,76). Eine isolierte Verbesserung der Grenzkontrollkapazität
reduzierte die Warteschlange dort um 94 Prozent, verlagerte den Engpass jedoch
stromabwärts und erhöhte die Gepäckwarteschlange von 120 auf 198 Passagiere,
ohne die Verweildauer zu verbessern.

Die integrierte Optimierung erreichte mit 78 Minuten die kürzeste
Verweildauer (minus 24 Prozent gegenüber dem Ausgangszustand). Bei einer
Nachfragesteigerung um 50 Prozent erzielte eine warteschlangengesteuerte
dynamische Intervention identische Ergebnisse wie die entsprechende statische
Konfiguration, die auf eine betriebliche Sättigungsgrenze hindeutet. Drei
unabhängig berechnete Verfahren - modellbasierte Merkmalsrelevanz,
SHAP-Attribution und die Power-BI-Key-Influencers-Analyse - weisen
übereinstimmend die Tageszeit als dominanten Einflussfaktor des Staurisikos
aus.

Der wissenschaftliche Beitrag besteht in einem integrierten,
ankunftszentrierten Framework, das eine proxybasierte Nachfrageschätzung mit
simulierten betrieblichen Wirkungen verknüpft und beides in eine
managementorientierte Entscheidungsunterstützung für einen westafrikanischen
Drehkreuzflughafen überführt - einen Kontext, der in der untersuchten
Literatur bislang nicht vertreten ist. Ergänzend wird ein methodischer Befund
berichtet, der die Nachfrageebene des Frameworks selbst einschränkt: Da sowohl
das Passagier- als auch das Gepäckaufkommen aus dem Flugzeugtyp abgeleitet
werden, ist die Gepäckbelastung eine deterministische lineare Transformation
der Passagierbelastung und keine eigenständige Prognosegröße; zudem nimmt der
Passagier-Proxy nimmt lediglich 10 verschiedene Werte an. Dies begrenzt, welche
Aussagen flugzeugtypbasierte Proxys über die ankunftsseitige Nachfrage
zulassen.

Wesentliche Einschränkungen sind entsprechend der Schätzung des Passagier-
und Gepäckaufkommens anhand des Flugzeugtyps anstelle erhobener Zählwerte, die
moderate Trefferquote des eingesetzten Klassifikators sowie die Beschränkung
auf einmalige Simulationsläufe ohne Wiederholungen.

\vspace{0.5cm}

\noindent
\textbf{Schlagwörter:}
prädiktive operative Intelligenz;
Ankunftsmanagement am Flughafen;
Passagierstromprognose;
Gepäckbelastung;
Stauklassifikation;
ereignisdiskrete Simulation;
Entscheidungsunterstützungssystem;
Internationaler Flughafen Blaise Diagne.